% Part: turing-machines
% Chapter: undecidability
% Section: unsolvability-decision-problem

\documentclass[../../../include/open-logic-section]{subfiles}

\begin{document}

\olfileid{tur}{und}{uns}
\olsection{مسئلهٔ تصمیم حل‌ناپذیر است}

\begin{thm}
\ollabel{thm:decision-prob}
مسئلهٔ تصمیم حل‌ناپذیر است: هیچ ماشین تورینگی چون~$D$ وجود ندارد که
هرگاه با نواری شامل !!a{sentence}ای از منطق مرتبه‌اول چون~$!B$ به‌عنوان
ورودی آغاز شود، سرانجام متوقف شود و اگر و تنها اگر $!B$ معتبر است
خروجی~$1$ و در غیر این صورت خروجی~$0$ بدهد.
\end{thm}

\begin{proof}
فرض کنید مسئلهٔ تصمیم حل‌پذیر باشد؛ یعنی ماشین تورینگی چون~$D$ وجود
داشته باشد. آنگاه می‌توانستیم مسئلهٔ توقف را چنین حل کنیم. ماشین
تورینگی چون~$E$ می‌سازیم که با دریافت شمارهٔ~$e$ از ماشین تورینگ
$M_e$ و ورودی~$w$، !!{sentence}ٔ متناظر
$!T(M_e,w)\lif !E(M_e,w)$ را محاسبه کند و در حالی متوقف شود که
چپ‌ترین خانهٔ نوار را پویش می‌کند. آنگاه ماشین~$E\concat D$ با دریافت
$e$ و~$w$ نخست $!T(M_e,w)\lif !E(M_e,w)$ را محاسبه و سپس ماشین
مسئلهٔ تصمیم~$D$ را روی آن ورودی اجرا می‌کرد. $D$ اگر و تنها اگر
$!T(M_e,w)\lif !E(M_e,w)$ معتبر باشد با خروجی~$1$، و در غیر این صورت
با خروجی~$0$ متوقف می‌شود. بنا بر
\olref[ver]{lem:halt-if-valid} و~\olref[ver]{lem:valid-if-halt}،
$!T(M_e,w)\lif !E(M_e,w)$ معتبر است اگر و تنها اگر $M_e$ روی
ورودی~$w$ متوقف شود. پس $E\concat D$ با ورودی~$e$ و~$w$ اگر و تنها
اگر $M_e$ روی~$w$ متوقف شود با خروجی~$1$، و در غیر این صورت با
خروجی~$0$ متوقف می‌شد. به بیان دیگر، $E\concat D$ مسئلهٔ توقف را حل
می‌کرد. اما از \olref[hal]{thm:halting-problem} می‌دانیم چنین ماشین
تورینگی نمی‌تواند وجود داشته باشد.
\end{proof}

\begin{cor}\ollabel{cor:undecidable-sat}%
تصمیم دربارهٔ ارضاپذیری یک !!{sentence}ٔ دلخواه از منطق مرتبه‌اول
امکان‌پذیر نیست.
\end{cor}

\begin{proof}
  فرض کنید ارضاپذیری با ماشین تورینگی چون~$S$ تصمیم‌پذیر باشد. آنگاه
  می‌توانستیم مسئلهٔ تصمیم را چنین حل کنیم: !!a{sentence}ٔ~$!B$ را به‌عنوان
  ورودی بگیرید و آن را یک خانه به راست ببرید. به خانهٔ~$1$ بازگردید
  و نماد~$\lnot$ را بنویسید.

  اکنون ماشین تورینگ~$S$ را اجرا کنید. این ماشین سرانجام با خروجی
  $1$ (اگر $\lnot !B$ ارضاپذیر باشد) یا~$0$ (اگر $\lnot !B$
  ارضاناپذیر باشد) روی نوار متوقف می‌شود. اگر در خانهٔ~$1$ یک
  $\TMstroke$ است آن را پاک کنید؛ و اگر خانهٔ~$1$ خالی است یک
  $\TMstroke$ بنویسید و سپس متوقف شوید.

  این ماشین تورینگ همیشه متوقف می‌شود و خروجی‌اش~$1$ است اگر و تنها
  اگر $\lnot !B$ ارضاناپذیر باشد، و در غیر این صورت~$0$. چون $!B$
  معتبر است اگر و تنها اگر $\lnot !B$ ارضاناپذیر باشد، ماشین اگر و
  تنها اگر $!B$ معتبر است خروجی~$1$ و در غیر این صورت خروجی~$0$
  می‌دهد؛ یعنی مسئلهٔ تصمیم را حل می‌کرد.
\end{proof}

\begin{explain}
پس هیچ ماشین تورینگی وجود ندارد که همیشه به پرسش «آیا $!B$ یک
!!{sentence}ٔ معتبر از منطق مرتبه‌اول است؟» پاسخ درست «بله» یا «خیر»
بدهد. بااین‌حال ماشین تورینگی \emph{وجود دارد} که همواره پاسخ درست
«بله» می‌دهد، اما اگر پاسخ «خیر» باشد اصلاً متوقف نمی‌شود. این از
قضیهٔ صحت و تمامیت منطق مرتبه‌اول و این واقعیت نتیجه می‌شود که
!!{derivation}s را می‌توان به‌طور مؤثر برشمرد.
\end{explain}

\begin{thm}
  \ollabel{thm:valid-ce}%
  اعتبار !!{sentence}sی منطق مرتبه‌اول نیمه‌تصمیم‌پذیر است: ماشین
  تورینگی چون~$E$ وجود دارد که هرگاه روی نواری شامل !!a{sentence}ای از
  منطق مرتبه‌اول چون~$!B$ به‌عنوان ورودی آغاز شود، اگر و تنها اگر
  $!B$ معتبر باشد سرانجام متوقف می‌شود و خروجی~$1$ می‌دهد، و در غیر
  این صورت متوقف نمی‌شود.
\end{thm}

\begin{proof}
  الگوریتمی مؤثر می‌تواند همهٔ !!{derivation}sی ممکنِ منطق مرتبه‌اول را
  یکی پس از دیگری تولید کند. ماشین~$E$ این کار را انجام می‌دهد و
  هنگامی که !!a{derivation}ی بیابد که $\Proves !B$ را نشان دهد، با
  خروجی~$1$ متوقف می‌شود. بنا بر قضیهٔ صحت، اگر $E$ با خروجی~$1$
  متوقف شود، دلیلش این است که~$\Entails !B$. بنا بر قضیهٔ تمامیت، اگر
  $\Entails !B$، !!a{derivation}ی وجود دارد که~$\Proves !B$ را نشان می‌دهد.
  چون $E$ همهٔ !!{derivation}sی ممکن را به‌طور منظم تولید می‌کند،
  سرانجام اشتقاقی را خواهد یافت که~$\Proves !B$ را نشان می‌دهد و پس
  با خروجی~$1$ متوقف خواهد شد.
\end{proof}

\end{document}
